\documentclass[11pt,a4paper]{article}
\usepackage[utf8]{inputenc}
\usepackage[T1]{fontenc}
\usepackage[margin=0.75in]{geometry}
\usepackage{enumitem}
\usepackage{titlesec}
\usepackage{xcolor}
\usepackage{hyperref}

\definecolor{primarycolor}{RGB}{33,37,41}
\definecolor{secondarycolor}{RGB}{0,102,204}
\definecolor{accentcolor}{RGB}{108,117,125}
\hypersetup{colorlinks=true,linkcolor=secondarycolor,urlcolor=secondarycolor,citecolor=secondarycolor}
\titleformat{\section}{\large\bfseries\color{primarycolor}}{}{0em}{}
\titlespacing*{\section}{0pt}{1em}{0.4em}
\setlist[itemize]{leftmargin=*,itemsep=0.15em,topsep=0.2em}
\pagestyle{empty}

\begin{document}
\begin{center}
{\Huge\bfseries\color{primarycolor} KATARZYNA WIŚNIEWSKA}\\[0.3em]
{\large\color{accentcolor} Pflegehelferin}\\[0.3em]
{\small Habsburgerstraße 41, 79104 Freiburg im Breisgau | k.wisniewska@example.com | +49 761 555 0412}
\end{center}

\section*{PROFIL}
Pflegehelferin mit Erfahrung in der stationären Altenpflege und fünf Jahren Berufserfahrung als Krankenschwester in Polen. Ruhig, zuverlässig und geduldig im Umgang mit älteren Menschen und Menschen mit Demenz. Erfahrung in der Grundpflege, Behandlungspflege nach LG1 und Pflegedokumentation. Deutschkenntnisse auf B2-Niveau.

\section*{BERUFSERFAHRUNG}
\textbf{Pflegehelferin, Seniorenzentrum Am Lindenhof, Freiburg} \hfill \textit{04/2020 – heute}
\begin{itemize}
\item Grundpflege für 10 bis 12 Bewohnerinnen und Bewohner pro Schicht: Körperpflege, Mobilisation, Lagerung und Hilfe beim Essen.
\item Behandlungspflege nach LG1: Blutdruck- und Blutzuckermessung, Kompressionsstrümpfe und Medikamente nach Plan reichen.
\item Pflegedokumentation in Vivendi PD.
\item Einsätze im Früh- und Spätdienst sowie an Wochenenden.
\end{itemize}

\textbf{Aushilfe Verkauf, Bäckerei Schwarzwaldkorn, Freiburg} \hfill \textit{09/2019 – 03/2020}
\begin{itemize}
\item Verkauf, Kasse und Reinigung.
\end{itemize}

\textbf{Krankenschwester (Pielęgniarka), Wojewódzki Szpital Specjalistyczny, Opole, Polen} \hfill \textit{10/2014 – 06/2019}
\begin{itemize}
\item Tätigkeit auf einer Station für Innere Medizin mit 32 Betten.
\item Grund- und Behandlungspflege, Infusionen, Wundversorgung und Vorbereitung von Visiten.
\end{itemize}

\section*{AUSBILDUNG}
Bachelor Pielęgniarstwo (Krankenpflege), Medizinische Universität Opole, 2014\\
Anerkennungsverfahren als Pflegefachfrau beim Regierungspräsidium Stuttgart: läuft

\section*{WEITERBILDUNG}
Behandlungspflege LG1 (Basiskurs), 2021\\
Expertenstandard Sturzprophylaxe (interne Schulung), 2023\\
Erste-Hilfe-Kurs, 2024

\section*{SPRACHEN}
Polnisch (Muttersprache)\\
Deutsch (telc Deutsch B1-B2 Pflege, B2)\\
Englisch (Grundkenntnisse)
\end{document}